\documentclass[conference]{IEEEtran}
\IEEEoverridecommandlockouts
% The preceding line is only needed to identify funding in the first footnote. If that is unneeded, please comment it out.
%Template version as of 6/27/2024

\usepackage{cite}
\usepackage{amsmath,amssymb,amsfonts}
\usepackage{algorithmic}
\usepackage{graphicx}
\usepackage{textcomp}
\usepackage{xcolor}
\def\BibTeX{{\rm B\kern-.05em{\sc i\kern-.025em b}\kern-.08em
    T\kern-.1667em\lower.7ex\hbox{E}\kern-.125emX}}
\begin{document}

\title{ Impacto de la cuantización en la reparación automática de vulnerabilidades mediante Code Language Models*\\

\thanks{Identify applicable funding agency here. If none, delete this.}
}

\author{\IEEEauthorblockN{1\textsuperscript{st} Bravo Arredondo Cristhian Matías}
\IEEEauthorblockA{\textit{Universidad Nacional de San Agustín} \\
Arequipa, Perú \\
cbravoa@unsa.edu.pe}
\and
\IEEEauthorblockN{2\textsuperscript{nd} Given Name Surname}
\IEEEauthorblockA{\textit{dept. name of organization (of Aff.)} \\
\textit{name of organization (of Aff.)}\\
City, Country \\
email address or ORCID}

}

\maketitle

\begin{abstract}
a
\end{abstract}

\begin{IEEEkeywords}
palabras clave
\end{IEEEkeywords}

\section{Introducción}

Según el NIST, el número de vulnerabilidades reportadas al sistema CVE creció un 263\,\% entre 2020 y 2025~\cite{nist2026nvd}. Corregirlas manualmente demanda mucho tiempo y esfuerzo~\cite{zhang2023pretrained}, lo que retrasa el desarrollo de los sistemas y los deja expuestos a ataques. Ante este escenario surge la reparación automática de programas, que busca generar parches con poca o ninguna intervención humana, y cuya viabilidad demostraron Weimer et al.~\cite{weimer2009genprog} con GenProg. Actualmente, esta tarea se aborda con \emph{Code Language Models} (CLM), modelos de lenguaje preentrenados principalmente con código fuente~\cite{jiang2023impact}.

En los últimos años, los CLM se han consolidado como una línea de investigación activa en seguridad de software: una revisión sistemática reciente reúne 58 estudios que los aplican a la detección y reparación de vulnerabilidades~\cite{zhou2025llmvuln}. No obstante, su efectividad en casos reales aún es limitada; en vulnerabilidades de Java, el mejor modelo evaluado corrigió apenas el 20,4\,\% de los casos, aunque el ajuste fino mejoró sus resultados~\cite{wu2023effective}. Siguiendo esta dirección, Bhandari et al.~\cite{bhandari2025patchlm} propusieron PatchLM, un CLM ajustado con correcciones de vulnerabilidades registradas como CVE, que superó a los modelos base CodeT5 y CodeLlama. Sin embargo, estos modelos demandan una gran cantidad de memoria: por ejemplo, DeepSeek-Coder de 33\,000 millones de parámetros requiere alrededor de 64\,GB de memoria de GPU~\cite{giagnorio2025quantizing}. Por ello se recurre a la cuantización, que representa los pesos del modelo con menos bits mediante un mapeo de números reales a enteros, tal como lo formalizaron Jacob et al.~\cite{jacob2018quantization}. En generación de código, Giagnorio et al.~\cite{giagnorio2025quantizing} mostraron que es posible bajar hasta 4 bits sin pérdida significativa y con cerca de 70\,\% menos de memoria, aunque con 3 y 2 bits el rendimiento cae de forma notable.

Sin embargo, el rendimiento de los CLM cuantizados en la reparación de programas ha sido poco evaluado, y los estudios disponibles se centran en errores funcionales~\cite{vallecillos2026smaller,kusama2025small}. Entre ellos, Vallecillos-Ruiz et al.~\cite{vallecillos2026smaller} observaron que los modelos cuantizados corrigen una cantidad similar de errores que los originales, pero no los mismos, y que la reducción de memoria, de hasta 85\,\%, se acompaña de un mayor tiempo de inferencia y consumo energético. Este cambio de comportamiento cobra relevancia en seguridad, donde se ha reportado que la cuantización a 4 bits incrementa la proporción de dependencias vulnerables recomendadas por el modelo~\cite{haque2025secure}. Pese a ello, en PatchLM~\cite{bhandari2025patchlm} el modelo CodeLlama se cuantizó a 4 bits por limitaciones de recursos, sin evaluar el efecto de esta decisión. Hasta donde sabemos, no existe un estudio que analice cómo la cuantización afecta la capacidad de los CLM para eliminar vulnerabilidades, la posible introducción de vulnerabilidades nuevas y el costo computacional de la reparación.

% --- Párrafo 4 (propósito, método y objetivo) y hoja de ruta: pendientes ---

\bibliographystyle{IEEEtran}
\bibliography{referencias}

\end{document}
